%---------------------------------------------------------------------
% Part four: the short-answer question. One combined question.
%
% Every paper's Q1 was ten short items, worth 20 marks together. Only
% terms this session could verify with confidence are included; a
% couple of narrower items from the scans (exact figures on minor,
% loosely-documented details) were left out rather than guessed --
% stated on the title page.
%---------------------------------------------------------------------

\section*{حصہ چہارم --- مختصر سوالات}
\markboth{مختصر سوالات}{}

\noindent
ہر پرچے کا \textbf{پہلا سوال لازمی} تھا: دس مختصر سوالات، کل بیس نمبر ---
یعنی فی سوال دو نمبر، وقت بیس منٹ سے زیادہ نہ لگائیے۔ ذیل میں وہ سوالات ہیں
جو واقعی پوچھے گئے:

\medskip
\begin{center}
\begin{tabular}{@{}p{3.1cm} p{11.3cm}@{}}
\toprule
\textbf{غزوہ اور سریہ کا فرق}
  & \textbf{غزوہ} وہ مہم جس میں نبی کریم \saw خود شریک ہوں؛ \textbf{سریہ}
    وہ جو کسی صحابیؓ کی قیادت میں بھیجی جائے اور نبی \saw خود شریک نہ ہوں۔
    {\small\color{bmgrey}[بہار \textenglish{2013}]} \\
\textbf{بیعتِ عقبہ}
  & یثرب کے باشندوں کی نبی کریم \saw سے دو بیعتیں --- عقبہ اولیٰ
    (\textenglish{621}ء، بارہ افراد، نیکی کی بیعت) اور عقبہ ثانیہ
    (\textenglish{622}ء، تہتر افراد، جان کی حفاظت کی بیعت)۔ تفصیل سوال
    \textenglish{6} میں۔ {\small\color{bmgrey}[بہار \textenglish{2013}]} \\
\textbf{تالیفِ قلوب}
  & زکوٰۃ کے مصارف میں سے ایک --- نئے مسلمانوں یا بااثر افراد کو ان کا دل
    اسلام کی طرف مائل رکھنے کے لیے دیا جانے والا حصہ۔
    {\small\color{bmgrey}[خزاں \textenglish{2012}]} \\
\textbf{حدیث اور سنت کا فرق}
  & \textbf{حدیث} نبی کریم \saw کا قول؛ \textbf{سنت} آپ \saw کا مستقل عمل
    و طرزِ زندگی --- سنت کا دائرہ حدیث سے وسیع ہے۔
    {\small\color{bmgrey}[خزاں \textenglish{2012}]} \\
\textbf{اصحابِ الفیل}
  & ابرہہ کی سرکردگی میں خانہ کعبہ پر ہاتھیوں کی فوج کے ساتھ حملہ آور
    لشکر (سنِ فیل، نبی کریم \saw کی پیدائش کا سال)۔ اللہ نے ابابیل پرندوں
    سے انہیں تباہ کیا (سورۃ الفیل)۔
    {\small\color{bmgrey}[بہار \textenglish{2013}]} \\
\textbf{رہبانیت}
  & دنیا و نکاح ترک کر کے مکمل گوشہ نشینی اختیار کرنا --- اسلام نے اس کی
    نفی کی: ''ہمارے دین میں رہبانیت نہیں'' (بیہقی)۔
    {\small\color{bmgrey}[بہار \textenglish{2013}]} \\
\textbf{حلفِ فضول}
  & زمانۂ جاہلیت میں قریش کے چند خاندانوں کا معاہدہ کہ مکہ میں کسی مظلوم
    کی مدد کریں گے۔ نبی کریم \saw نے جوانی میں اس میں شرکت کی اور نبوت کے
    بعد بھی فرمایا کہ آج بھی بلایا جاؤں تو شریک ہوں۔
    {\small\color{bmgrey}[خزاں \textenglish{2017}]} \\
\textbf{فترتِ وحی}
  & پہلی وحی (اقرأ) کے نزول کے بعد وحی کے سلسلے میں کچھ عرصے کا وقفہ، جو
    سورۃ المدثر کے نزول پر ختم ہوا۔ تفصیل سوال \textenglish{3} میں۔
    {\small\color{bmgrey}[خزاں \textenglish{2017}]} \\
\textbf{لا وصیۃ لوارث}
  & شرعی وارث کے لیے مقررہ حصے کے علاوہ الگ سے وصیت جائز نہیں --- حدیث:
    ''کسی وارث کے لیے وصیت نہیں'' (ترمذی)۔
    {\small\color{bmgrey}[بہار \textenglish{2013}]} \\
\textbf{اجود من الریح المرسلۃ}
  & ''بھیجی ہوئی ہوا سے بھی زیادہ سخی'' --- نبی کریم \saw کی سخاوت کی
    تشبیہ، خصوصاً رمضان میں (بخاری)۔ تفصیل سوال \textenglish{16} میں۔
    {\small\color{bmgrey}[خزاں \textenglish{2017}]} \\
\textbf{ایثار اور اسراف کا فرق}
  & \textbf{ایثار} اپنی ضرورت پر دوسرے کو ترجیح دینا (پسندیدہ)؛
    \textbf{اسراف} حد سے تجاوز کر کے فضول خرچی کرنا (ناپسندیدہ)۔
    {\small\color{bmgrey}[خزاں \textenglish{2017}]} \\
\bottomrule
\end{tabular}
\end{center}
